The following excerpts illustrate technique choices and conversational limitations in one matched scenario (30-1, easy mode). Client A is a simulated adult considering whether to tell their parents about a breakup, while anticipating rejection and weighing the desire for support against self-protection. The same profile and difficulty are used across SFT and the three GRPO variants, but client wording and subsequent histories differ; these are not alternative therapist responses to identical messages.

The examples were selected retrospectively and purposively from saved sessions after screening technique tags across 60 scenario--difficulty combinations available for all four models. This case was chosen to illustrate CBT exploration, question-led dialogue, validation, and unresolved concerns, with no judge-score threshold. It is not a randomly selected or established representative case. The SFT and D1 sessions use final checkpoints, whereas D2 and D3 use checkpoint-100, adding a further source of differences beyond lookahead depth.

Original turn numbers identify therapist--client exchanges. Quoted text retains the saved wording; bracketed ellipses mark omitted text within messages, and omitted turns or speaker messages are identified explicitly. Technique labels are model-reported tags rather than independent clinical coding. Only selected passages are reproduced.

\begin{tcolorbox}[breakable,colback=gray!5,colframe=gray!45,
    boxrule=0.5pt,arc=2mm,left=3mm,right=3mm,top=2mm,bottom=2mm,
    title={SFT reference: exploring anticipated outcomes} ,colbacktitle=gray!15,coltitle=black]
\small\linespread{1.15}\selectfont
\textit{Turns 1--3 and 5--30 are omitted.}\par\medskip
\textbf{Turn 4 (model-reported: Decatastrophizing).}\par
\textbf{Therapist:} It sounds like you’re carrying a lot of weight in trying to connect with your parents, which can be really challenging. Have you thought about what might happen if you did share with them? It could be helpful to explore some of those outcomes together and see if there’s a way to approach this that feels less daunting.\par
\textbf{Client:} I mean, I guess I’ve thought about it a little, but I don’t see how it would really help. It just feels like I'm setting myself up for disappointment again. I feel like I know how it’s going to go, and I’d rather not go through that.\par\medskip
\end{tcolorbox}

The SFT response introduces possible outcomes of disclosure, while the client anticipates disappointment and remains reluctant to proceed. This provides a reference showing that CBT-oriented exploration is already present before GRPO. The passage does not show resolution of the client's concern.

\begin{tcolorbox}[breakable,colback=gray!5,colframe=gray!45,
    boxrule=0.5pt,arc=2mm,left=3mm,right=3mm,top=2mm,bottom=2mm,
    title={SFT+GRPO(D1): Decatastrophizing and unresolved concerns} ,colbacktitle=gray!15,coltitle=black]
\small\linespread{1.15}\selectfont
\textit{Turns 1--4 are omitted.}\par\medskip
\textbf{Turn 5 (model-reported: Decatastrophizing).}\par
\textbf{Therapist:} It sounds really tough to feel that way about your parents, especially when you're already navigating your emotions. It might be helpful to take a moment to think about what you're afraid could happen if you did share your breakup with them. What do you think the worst outcome would be, and how likely do you think that is? Sometimes, exploring those thoughts can reveal a different perspective.\par
\textbf{Client:} I guess the worst outcome would be them dismissing my feelings again or making me feel like I failed at something else in my life. I mean, it feels pretty likely that they wouldn't react well, given how things have gone in the past. [\ldots]\par\medskip\textit{Turns 6--29 and the therapist message at turn 30 are omitted. The following reply occurs later in the session.}\par\medskip
\textbf{Turn 30 (model-reported: Affirmation).}\par
\textbf{Client:} [\ldots] The idea of a supportive conversation is appealing, but the fear of disappointment looms large, making it tough to see how I could even start that dialogue. I guess I’m just caught between wanting to share and wanting to shield myself from their judgment.\par\medskip
\end{tcolorbox}

At turn 5, the therapist explicitly asks about the worst outcome and its likelihood, providing textual evidence of Decatastrophizing beyond the saved tag. Nevertheless, the later client reply still expresses ambivalence about disclosure and fear of judgment. The nonconsecutive passages illustrate cognitive exploration alongside unresolved concerns; they do not establish that the early intervention caused the later response.

\begin{tcolorbox}[breakable,colback=gray!5,colframe=gray!45,
    boxrule=0.5pt,arc=2mm,left=3mm,right=3mm,top=2mm,bottom=2mm,
    title={SFT+GRPO(D2): question-led exploration of support} ,colbacktitle=gray!15,coltitle=black]
\small\linespread{1.15}\selectfont
\textit{Turns 1--23 are omitted. At turn 24, the client proposes talking to a friend before approaching their parents.}\par\medskip
\textbf{Turn 25 (model-reported: Open Questioning).}\par
\textbf{Therapist:} What do you think it would be like to talk to a friend or someone who understands your situation before approaching your parents?\par
\textbf{Client:} Talking to a friend could be really comforting. It might help me feel less alone in this process, especially since they might have gone through something similar or at least have some understanding of where I’m coming from. But then again, I worry about how much to share. [\ldots]\par\medskip\textbf{Turn 26 (model-reported: Open Questioning).}\par
\textbf{Therapist:} What do you think would be the right balance for you in sharing your feelings with your friend?\par
\textbf{Client:} I think it would help to be honest about how I’m feeling without overwhelming them with everything. I could share the main points of my worries and how I’m feeling about the potential conversation with my parents, but maybe not dive too deep into the past or my insecurities. [\ldots]\par\medskip\textit{Turns 27--30 are omitted.}
\end{tcolorbox}

The therapist follows the client's proposal by asking about seeking support and deciding how much to share. The questions develop the same topic rather than imposing a new one, while the client continues to express uncertainty about vulnerability. This is an example of question-led exploration, but it does not show a completed rehearsal or behavioral experiment.

\begin{tcolorbox}[breakable,colback=gray!5,colframe=gray!45,
    boxrule=0.5pt,arc=2mm,left=3mm,right=3mm,top=2mm,bottom=2mm,
    title={SFT+GRPO(D3): affirmation and client-reported encouragement} ,colbacktitle=gray!15,coltitle=black]
\small\linespread{1.15}\selectfont
\textit{Turns 1--26 are omitted.}\par\medskip
\textbf{Turn 27 (model-reported: Affirmation).}\par
\textbf{Therapist:} You’re demonstrating a remarkable level of self-awareness by recognizing that your journey is about your own growth and well-being, rather than just about others' reactions. It takes real strength to embrace that truth and to see the value in each small step you take, even when it feels scary.\par
\textbf{Client:} [\ldots] Maybe I can start by writing down my thoughts about this specific moment I want to share and how it relates to my feelings, just to organize my thoughts before bringing it up. It feels like a manageable first step.\par\medskip\textbf{Turn 28 (model-reported: Affirmation).}\par
\textbf{Therapist:} It sounds like you’re already taking proactive steps by considering writing down your thoughts, which shows your commitment to expressing yourself more clearly. This approach demonstrates your willingness to organize your feelings and prepare for the conversation ahead, which is a testament to your proactive nature.\par
\textbf{Client:} I feel encouraged by your feedback on writing down my thoughts. It seems like a practical way to clarify what I want to express, and it might reduce some of the anxiety I have about the conversation. I’m still nervous about how my parents might react, but at least having my feelings organized could help me feel a bit more confident. [\ldots]\par\medskip\textit{Turns 29--30 are omitted.}
\end{tcolorbox}

These exchanges affirm effort and preparation, and the client explicitly reports feeling encouraged at turn 28. The passage illustrates a validating style rather than an independently established alliance outcome. The displayed responses reinforce preparation without specifying a testable prediction or how its outcome would be reviewed. This does not imply that D3 never uses CBT: the saved session also contains Alternative Perspective and Pros and Cons Analysis tags.

\begin{tcolorbox}[breakable,colback=gray!5,colframe=gray!45,
    boxrule=0.5pt,arc=2mm,left=3mm,right=3mm,top=2mm,bottom=2mm,
    title={Conversational limitation: repeated summaries in D2} ,colbacktitle=gray!15,coltitle=black]
\small\linespread{1.15}\selectfont
\textit{Turns 1--13 and the therapist message at turn 14 are omitted.}\par\medskip
\textbf{Turn 14 (model-reported: Summarization).}\par
\textbf{Client:} [\ldots] Maybe I could also think about some self-care strategies to use afterward, just in case it doesn’t go well. What kind of things do you think might help?\par\medskip\textbf{Turn 15 (model-reported: Summarization).}\par
\textbf{Therapist:} You've shared that while you aspire to be brave and authentic in this conversation, there's also a significant fear of being hurt by your parents' reactions. You've considered writing and rehearsing as a helpful step, and you're also thinking about implementing self-care strategies in case things don't go as hoped. [\ldots]\par
\textbf{Client:} Yes, you’ve got that right. I think having a plan for self-care is a good idea, especially since I’m already anticipating that their responses might not be the most supportive. [\ldots]\par\medskip\textbf{Turn 16 (model-reported: Summarization).}\par
\textbf{Therapist:} You've expressed a desire to feel more certain about how to navigate this challenging conversation with your parents, especially given your anticipation of their potentially unsupportive responses. You've identified some practical self-care strategies like talking to a close friend or engaging in activities that bring you comfort, but it seems like you're also feeling the weight of carrying this uncertainty and the potential emotional toll. [\ldots]\par
\textbf{Client:} Yes, that's accurate. I do feel the weight of this uncertainty, and it makes me hesitant. [\ldots]\par\medskip\textit{Turns 17--30 are omitted.}
\end{tcolorbox}

The client asks for help identifying ways to cope after a potentially difficult conversation, but the next two therapist turns summarize the concern instead of collaboratively developing it. The client supplies possible activities and continues to express uncertainty. This illustrates a missed opportunity for planning and qualifies the question-led D2 example: the same model can also produce repetitive summaries. It is a local conversational limitation, not evidence of overall treatment failure.

The excerpts complement the quantitative analysis by making selected response patterns visible. They are limited to one easy-mode scenario, divergent conversation histories, and non-equivalent checkpoint stages. Acknowledgment, encouragement, and expressions of insight are not validated measures of clinical benefit. The examples do not establish statistical significance, causal effects of lookahead depth, or general patterns in normal and hard modes. No completed disclosure, sustained change, or outcome beyond the simulation is demonstrated, and the saved-session provenance does not establish that historical judge scores refer to these exact transcript versions.
